\section{Demultiplexador (DEMUX 1x4 com portas NOR)}

\subsection{Tabela Verdade e Expressões Booleanas}
Abaixo apresentamos a tabela verdade para o DEMUX 1x4. Considerando $En$ como a entrada de habilitação, $S_1$ e $S_0$ como chaves de seleção e $D$ como o dado de entrada.

\begin{table}[H]
    \centering
    \caption{Tabela Verdade do DEMUX 1x4}
    \begin{tabular}{cccc|cccc}
        \toprule
        \textbf{$En$} & \textbf{$S_1$} & \textbf{$S_0$} & \textbf{$D$} & \textbf{$Y_0$} & \textbf{$Y_1$} & \textbf{$Y_2$} & \textbf{$Y_3$} \\
        \midrule
        0             & X              & X              & X            & 0              & 0              & 0              & 0              \\
        1             & 0              & 0              & D            & D              & 0              & 0              & 0              \\
        1             & 0              & 1              & D            & 0              & D              & 0              & 0              \\
        1             & 1              & 0              & D            & 0              & 0              & D              & 0              \\
        1             & 1              & 1              & D            & 0              & 0              & 0              & D              \\
        \bottomrule
    \end{tabular}
\end{table}

As expressões booleanas na forma de Soma de Produtos (SOP) seriam:
\begin{align*}
    Y_0 & = En \cdot \overline{S_1} \cdot \overline{S_0} \cdot D \\
    Y_1 & = En \cdot \overline{S_1} \cdot S_0 \cdot D            \\
    Y_2 & = En \cdot S_1 \cdot \overline{S_0} \cdot D            \\
    Y_3 & = En \cdot S_1 \cdot S_0 \cdot D
\end{align*}

Para a implementação solicitada (apenas portas NOR), aplicamos o Teorema de De Morgan ($A \cdot B = \overline{\overline{A} + \overline{B}}$):
% Substitua pelas equações finais manipuladas
\begin{align*}
    Y_0 & = \overline{ \overline{En} + S_1 + S_0 + \overline{D} }                       \\
    Y_0 & = \overline{ \overline{En} + S_1 + S_0 + \overline{D} }                       \\
    Y_1 & = \overline{ \overline{En} + S_1 + \overline{S_0} + \overline{D} }            \\
    Y_2 & = \overline{ \overline{En} + \overline{S_1} + S_0 + \overline{D} }            \\
    Y_3 & = \overline{ \overline{En} + \overline{S_1} + \overline{S_0} + \overline{D} }
    % Adicione as demais aqui...
\end{align*}

\subsection{Diagrama Esquemático}
Para a obtenção do esquemático, as expressões originais foram convertidas para a base NOR invertendo duplamente a saída e aplicando De Morgan nas portas AND originais.

\begin{figure}[H]
    \centering
    % Substitua "nome_da_imagem.png" pelo arquivo do seu esquemático
    % \includegraphics[width=0.8\textwidth]{esquematico_4_1.png}
    \caption{Diagrama esquemático do DEMUX 1x4 (Portas NOR).}
\end{figure}

\subsection{Diagrama de Temporização}
\begin{figure}[H]
    \centering
    % \includegraphics[width=0.9\textwidth]{tempo_4_1.png}
    \caption{Simulação funcional (diagrama de temporização) do DEMUX 1x4.}
\end{figure}